\section{Research Questions}
\label{sec:rqs}

Fine-tuning a code model on obfuscated programs raises its accuracy on those programs by roughly
nineteen points. Every question below asks what that number is made of. Adaptation can raise forward
accuracy by teaching the model to \emph{read} the program better, or by teaching it the surface a
particular transform leaves behind; only the first deserves to be called robustness. We separate them
with two tests that share no machinery, committed before any method was on the table:
\begin{itemize}[leftmargin=*, topsep=2pt, itemsep=1pt]
  \item \textbf{Composition.} Does competence on single transforms survive being stacked with a
        transform family the model has never seen?
  \item \textbf{Reversal.} Does it survive the task being run backwards --- given a return value,
        produce a call that yields it? No adapter is trained on this direction, so any backward gain
        is transferred program understanding rather than task fitting.
\end{itemize}

\noindent\textbf{RQ1: Do standard adaptation methods generalise?} We apply both tests to the field's
existing strategies --- inference-time routing over per-transform specialists, weight-space merging of
those specialists, and broad multi-transform fine-tuning --- treated as baselines rather than
proposals, and then locate what they learn instead, because that is what determines the remedy.

\noindent\textbf{RQ2: Does anchoring to clean-code behaviour repair what adaptation breaks?} The
diagnosis names its own remedy: anchor the model to its own behaviour on the unobfuscated parent of
each program. We derive a paired-consistency objective, isolate its mechanism by ablating the teacher,
and measure it across eight models and four lineages.

\noindent\textbf{RQ3: Does the anchored model survive the same tests?} We re-run RQ1's two tests on
the anchored model and state the bound: a new capability on unseen obfuscation, or a preservation
mechanism that protects the teacher's competence from surface-level change.
